\documentclass[11pt]{article}
\usepackage[T1]{fontenc}
\usepackage[utf8]{inputenc}
\usepackage[spanish]{babel}
\usepackage{amsmath,amssymb}
\usepackage[a4paper,margin=2.3cm]{geometry}
\usepackage{graphicx}
\usepackage{tikz}
\usetikzlibrary{arrows.meta,babel}
\IfFileExists{separacion_variables_visualizacion.pdf}{%
  \newcommand{\includevisualizacion}[1][]{\includegraphics[##1]{separacion_variables_visualizacion.pdf}}%
}{%
  \usepackage[inkscapelatex=false]{svg}
  \newcommand{\includevisualizacion}[1][]{\includesvg[##1]{separacion_variables_visualizacion}}%
}
\spanishdecimal{.}
\setlength{\parindent}{0pt}
\setlength{\parskip}{0.35em}
\setlength{\abovedisplayskip}{4pt}
\setlength{\belowdisplayskip}{4pt}
\setlength{\abovedisplayshortskip}{4pt}
\setlength{\belowdisplayshortskip}{4pt}

\newcommand{\encabezado}[3]{%
  {\Large\bfseries #1\par}
  \vspace{2pt}
  {\normalsize #2\hfill #3\par}
  \vspace{6pt}\hrule\vspace{8pt}}

\newcommand{\problema}[2]{%
  \vspace{0.55em}
  \noindent\textbf{#1.}\enspace{\small #2}\par
  \vspace{0.3em}}

\begin{document}
\encabezado{Potencial y campo electrostático por separación de variables en un cilindro}
{Electrodinámica I}{P.S., J.A., J.R.}

\problema{1}{Una superficie cilíndrica de radio $R$, infinita en la
dirección $z$, separa dos regiones de vacío sin carga volumétrica. El potencial
no depende de $z$ y está impuesto sobre la superficie:
\begin{equation}
  V(R,\phi)=V_0\sin(2\phi).
  \label{eq:borde}
\end{equation}
El potencial es finito en el eje y tiende a cero cuando $s\to\infty$.
Determine $V$ y $\mathbf E$ en ambas regiones y la densidad de carga
superficial necesaria para sostener la condición de borde. La superficie se
trata como una interfaz con carga; no como un único conductor en equilibrio,
pues su potencial varía con $\phi$.}

\begin{figure}[htbp]
\centering
\begin{tikzpicture}[
  scale=1.05,>=Stealth,
  line cap=round,line join=round,
  every node/.style={font=\small}
]
  \def\rad{2.0}
  % El sombreado identifica solamente los sectores positivos.
  \fill[black!10] (0,0)--(0:\rad) arc (0:90:\rad)--cycle;
  \fill[black!10] (0,0)--(180:\rad) arc (180:270:\rad)--cycle;
  \draw[->,semithick] (-2.55,0)--(2.70,0) node[right] {$x$};
  \draw[->,semithick] (0,-2.48)--(0,2.55) node[above] {$y$};
  \draw[thick] (0,0) circle (\rad);
  \fill (0,0) circle (1.2pt) node[below left] {$O$};

  \node at (60:1.35) {$+$};
  \node at (135:1.35) {$-$};
  \node at (225:1.35) {$+$};
  \node at (315:1.35) {$-$};

  \coordinate (P) at (35:\rad);
  \draw[->] (0,0)--(P) node[pos=0.69,above] {$R$};
  \draw[->] (0.48,0) arc (0:35:0.48);
  \node at (17:0.67) {$\phi$};
  \fill (P) circle (1.4pt);
  \draw[->,thick] (P)--++(35:0.65)
       node[above right] {$\hat{\mathbf s}$};
  \draw[->,thick] (P)--++(125:0.65)
       node[above left] {$\hat{\mathbf e}_{\phi}$};
\end{tikzpicture}
\caption{Corte transversal visto desde $+z$. El sombreado y los signos
indican $V(R,\phi)>0$ y $V(R,\phi)<0$ en sectores alternados. La normal
$\hat{\mathbf s}$ apunta del interior ($s<R$) al exterior ($s>R$).}
\label{fig:geometria}
\end{figure}

\textbf{Solución:}

\textbf{1. Separación de variables.}\par
En ausencia de carga volumétrica,
\begin{equation}
 \nabla^2V=\frac1s\frac{\partial}{\partial s}
 \left(s\frac{\partial V}{\partial s}\right)
 +\frac1{s^2}\frac{\partial^2V}{\partial\phi^2}=0.
 \label{eq:laplace}
\end{equation}
Con $V(s,\phi)=S(s)\Phi(\phi)$, multiplicamos por $s^2/(S\Phi)$:
\[
 \frac{s}{S}\frac{d}{ds}\left(s\frac{dS}{ds}\right)
 =-\frac{\Phi''}{\Phi}=m^2.
\]
Por lo tanto,
\begin{equation}
 \Phi''+m^2\Phi=0,\qquad s^2S''+sS'-m^2S=0.
 \label{eq:edos}
\end{equation}
La periodicidad $\Phi(\phi+2\pi)=\Phi(\phi)$ exige $m=0,1,2,\ldots$.
Para $m\geq1$ las soluciones son
\[
 \Phi_m=a_m\cos(m\phi)+b_m\sin(m\phi),\qquad
 S_m=c_ms^m+d_ms^{-m}.
\]
Para $m=0$, la solución angular periódica es constante y las dos soluciones
radiales son $1$ y $\ln s$. El logaritmo diverge en el eje y en el infinito;
la constante exterior tampoco cumple $V\to0$.

\textbf{2. Potenciales.}\par
La finitud en $s=0$ elimina $s^{-m}$ y $\ln s$
del interior. El decaimiento al infinito elimina $s^m$, $\ln s$ y la constante
del exterior. Por superposición,
\begin{align*}
 V_{\mathrm{int}}(s,\phi)
 &=a_0+\sum_{m=1}^{\infty}s^m
   [a_m\cos(m\phi)+b_m\sin(m\phi)],\\
 V_{\mathrm{ext}}(s,\phi)
 &=\sum_{m=1}^{\infty}s^{-m}
   [c_m\cos(m\phi)+d_m\sin(m\phi)].
\end{align*}
La condición \eqref{eq:borde} contiene solamente el armónico $\sin(2\phi)$.
La ortogonalidad de Fourier anula todos los demás coeficientes. Al imponer
el valor $V_0$ en $s=R$ queda
\begin{equation}
 \boxed{V_{\mathrm{int}}=V_0\left(\frac{s}{R}\right)^2\sin(2\phi),
 \qquad
 V_{\mathrm{ext}}=V_0\left(\frac{R}{s}\right)^2\sin(2\phi).}
 \label{eq:potenciales}
\end{equation}
Ambas expresiones valen $V_0\sin(2\phi)$ al llegar a la superficie.

\textbf{3. Campo eléctrico.}\par
Como no hay dependencia de $z$,
\[
 \mathbf E=-\nabla V
 =-\frac{\partial V}{\partial s}\hat{\mathbf s}
 -\frac1s\frac{\partial V}{\partial\phi}\hat{\mathbf e}_{\phi}.
\]
Al derivar \eqref{eq:potenciales},
\begin{align}
 \boxed{\mathbf E_{\mathrm{int}}
 =-\frac{2V_0s}{R^2}
 [\sin(2\phi)\hat{\mathbf s}
 +\cos(2\phi)\hat{\mathbf e}_{\phi}]}&\quad(s<R),
 \label{eq:eint}\\
 \boxed{\mathbf E_{\mathrm{ext}}
 =\frac{2V_0R^2}{s^3}
 [\sin(2\phi)\hat{\mathbf s}
 -\cos(2\phi)\hat{\mathbf e}_{\phi}]}&\quad(s>R).
 \label{eq:eext}
\end{align}
En particular, las componentes tangenciales coinciden en la frontera:
\[
 E_{\phi}^{\mathrm{int}}(R,\phi)
 =E_{\phi}^{\mathrm{ext}}(R,\phi)
 =-\frac{2V_0}{R}\cos(2\phi).
\]

\textbf{4. Carga superficial.}\par
Con la normal de la figura y vacío a ambos
lados, la condición de salto de Gauss es
\[
 \sigma(\phi)=\varepsilon_0
 [E_{s}^{\mathrm{ext}}(R,\phi)-E_{s}^{\mathrm{int}}(R,\phi)].
\]
Como $E_s^{\mathrm{ext}}=2V_0\sin(2\phi)/R$ y
$E_s^{\mathrm{int}}=-2V_0\sin(2\phi)/R$ en $s=R$,
\begin{equation}
 \boxed{\sigma(\phi)=\frac{4\varepsilon_0V_0}{R}\sin(2\phi).}
 \label{eq:sigma}
\end{equation}
La carga neta por unidad de longitud es nula:
$\int_0^{2\pi}\sigma(\phi)R\,d\phi=0$.

\clearpage
\section*{Visualización numérica}
La figura~\ref{fig:numerica} evalúa \eqref{eq:potenciales} y
\eqref{eq:eint}--\eqref{eq:sigma} para $R=V_0=1$. El mapa de potencial
muestra el modo angular de cuatro sectores; las flechas son valores locales del
campo, no líneas de campo. Los paneles inferiores comparan la condición de
borde $V(R,\phi)/V_0=\sin(2\phi)$ con la carga normalizada
$\sigma R/(4\varepsilon_0V_0)=\sin(2\phi)$.
\begin{figure}[htbp]
 \centering
 \includevisualizacion[width=0.98\textwidth]
 \caption{(a) Potencial adimensional. (b) Magnitud y dirección del campo
 eléctrico. (c) Potencial impuesto sobre la superficie. (d) Densidad
 superficial normalizada. La línea circular marca $s=R$.}
 \label{fig:numerica}
\end{figure}

\end{document}
